% ============================================================
% preamble.tex — Preámbulo compartido de SemArq1_TeamBales
% Los fragmentos (solucion.tex, notas.tex) NO cargan paquetes:
% todo lo que necesiten se carga aquí.
% ============================================================

% --- Página: márgenes de 2,5 cm (la portada alinea su texto con este margen) ---
\usepackage[margin=2.5cm]{geometry}

% --- Codificación y fuentes (ATS: tildes extraíbles con pdftotext) ---
\usepackage[T1]{fontenc}
\usepackage{lmodern}

% --- Idioma ---
% es-nolists: babel no redefine las listas (enumitem las controla).
% es-tabla:   «Tabla» en lugar de «Cuadro».
\usepackage[spanish, es-nolists, es-tabla]{babel}
\usepackage{csquotes}
\addto\captionsspanish{\renewcommand{\contentsname}{Índice}}

% --- Bibliografía: autor-año, una sola lista final ---
\usepackage[backend=biber, style=authoryear, natbib=true]{biblatex}
\addbibresource{core.bib}
% Fuentes nuevas de cada pregunta e investigación (solo si el archivo existe).
\newcommand{\bibsiexiste}[1]{\IfFileExists{#1}{\addbibresource{#1}}{}}
\bibsiexiste{preguntas/p1-trade-offs/referencias.bib}
\bibsiexiste{preguntas/p2-simetria-hexagono/referencias.bib}
\bibsiexiste{preguntas/p3-ddd-bounded-context/referencias.bib}
\bibsiexiste{preguntas/p4-inversion-dependencias/referencias.bib}
\bibsiexiste{preguntas/p5-kitty-problem/referencias.bib}
\bibsiexiste{preguntas/p6-ncapas-calidad/referencias.bib}
\bibsiexiste{investigacion/p1-regla-dependencia/fuentes.bib}
\bibsiexiste{investigacion/p2-hexagonal-onion/fuentes.bib}
\bibsiexiste{investigacion/p3-bounded-context/fuentes.bib}
\bibsiexiste{investigacion/p4-dip-interfaces/fuentes.bib}
\bibsiexiste{investigacion/p5-kitty-ocp/fuentes.bib}
\bibsiexiste{investigacion/p6-open-layers/fuentes.bib}

% --- Matemáticas, tablas, figuras y código ---
\usepackage{amsmath}
\usepackage{graphicx}
\usepackage{booktabs}
\usepackage{array}
\usepackage{tabularx}
\usepackage{float}
\usepackage{enumitem}
\usepackage{listings}
\usepackage{xcolor}
\usepackage{tikz}
% babel: evita el choque entre los atajos de babel-spanish (< y >) y las flechas de TikZ.
\usetikzlibrary{arrows.meta, positioning, shapes.geometric, calc, babel}

% --- Microtipografía: mejora el ajuste de las líneas (texto real, sin efecto en ATS) ---
\usepackage{microtype}

% --- Hipervínculos (al final, como recomienda hyperref) ---
\usepackage{hyperref}

% Estilos compartidos de diagramas (copiables por P3 y P4)
\tikzset{
  bloque/.style={rectangle, draw, thick, minimum width=2.4cm, minimum height=1.1cm, align=center, font=\sffamily\small},
  flecha/.style={-{Stealth[length=3mm]}, thick, shorten >=1pt, shorten <=1pt},
  capa/.style={rectangle, draw=gray!60, thick, rounded corners, minimum width=3.2cm, minimum height=1.6cm, align=center, font=\sffamily\footnotesize, fill=gray!8},
  contexto/.style={draw, thick, rounded corners=4pt, minimum width=4cm, minimum height=2.6cm, align=center, font=\sffamily\small, fill=blue!6},
  linea-integracion/.style={-{Stealth}, thick, dashed},
  acl/.style={rectangle, draw, thick, minimum width=2.6cm, minimum height=1cm, align=center, font=\sffamily\footnotesize, fill=orange!15},
  dependencia/.style={-{Latex[length=2.5mm]}, thick, shorten >=1pt, shorten <=1pt},
  modulo/.style={rectangle, draw, thick, minimum width=2.8cm, minimum height=0.9cm, align=center, font=\sffamily\small},
  interfaz/.style={rectangle, draw, thick, dashed, minimum width=2.8cm, minimum height=0.9cm, align=center, font=\sffamily\small\itshape, fill=yellow!12},
}

% Ajustes de composición
% - \emergencystretch: permite repartir más espacio en líneas con palabras que no se
%   pueden cortar (p. ej. nombres de clases en \texttt), en vez de salirse del margen.
% - biburl*penalty: permite cortar las URL largas de la bibliografía.
\setlength{\emergencystretch}{3em}
% - \tolerance: acepta líneas algo más espaciadas antes que dejar una que se salga del margen.
\tolerance=1000
\setcounter{biburlnumpenalty}{100}
\setcounter{biburlucpenalty}{100}
\setcounter{biburllcpenalty}{100}

% Capítulos seguidos, sin salto de página: cada capítulo (introducción, preguntas,
% conclusiones, bibliografía, uso de IA) empieza a continuación del anterior.
% Se quita el \clearpage de \chapter y se reduce el espacio antes del título.
\makeatletter
\patchcmd{\chapter}{\if@openright\cleardoublepage\else\clearpage\fi}{\par\bigskip}{}{}
\patchcmd{\@makeschapterhead}{\vspace*{50\p@}}{\vspace*{10\p@}}{}{}
% Títulos de capítulo: sin la línea «Capítulo N» (el título ya dice «P4: …»),
% en \LARGE en lugar de \Huge y con menos espacio. El número del capítulo se
% conserva internamente para numerar las secciones (4.1, 4.2…) y el índice.
\renewcommand{\@makechapterhead}[1]{%
  \vspace*{10\p@}%
  {\parindent \z@ \raggedright \normalfont
    \interlinepenalty\@M
    \LARGE\bfseries #1\par\nobreak
    \vskip 20\p@}}
\patchcmd{\@makeschapterhead}{\Huge \bfseries #1}{\LARGE\bfseries #1}{}{}
\patchcmd{\@makeschapterhead}{\vskip 40\p@}{\vskip 20\p@}{}{}
\makeatother

% Formato único de citas: (Autor, año, p. XX) y (Autor, año, pp. XX–YY).
% - Prefijos «p.» y «pp.» también cuando biblatex los añade solo (en español pondría
%   «pág.» y «págs.»).
% - \mknormrange: los rangos de páginas se escriben con guion largo (154–155).
\DefineBibliographyStrings{spanish}{page = {p\adddot}, pages = {pp\adddot}}
% biblatex en español separa los rangos con guion corto; aquí se usa el guion largo.
\DefineBibliographyExtras{spanish}{\protected\def\bibrangedash{\textendash\penalty\hyphenpenalty}}
\DeclareFieldFormat{postnote}{\mkpageprefix[pagination][\mknormrange]{#1}}
\DeclareFieldFormat{multipostnote}{\mkpageprefix[pagination][\mknormrange]{#1}}

% Wrappers de citas: biblatex elige «p.» (una página) o «pp.» (rango o lista).
\newcommand{\fuente}[2]{\citep[#2]{#1}}
\newcommand{\nombrefuente}[1]{\citet{#1}}

% Referencias bibliográficas en formato APA con viñetas
\defbibenvironment{bibliography}
  {\list
     {\textbullet}
     {\setlength{\labelwidth}{1em}%
      \setlength{\labelsep}{0.5em}%
      \setlength{\leftmargin}{1.5em}%
      \setlength{\itemindent}{0pt}%
      \setlength{\itemsep}{0.4em}%
      \setlength{\parsep}{0pt}}}
  {\endlist}
  {\item}

% Metadatos PDF (ATS)
\hypersetup{
  colorlinks=true,
  linkcolor=blue,
  citecolor=black,
  urlcolor=blue,
  pdftitle={Seminario de Arquitecturas de Software: Un Análisis Comparativo de Estilos Basados en Capas y Dominio},
  pdfauthor={Team Bales},
}
